\begin{center}
Problema E

Estória da Estônia

{\small Nome do arquivo: E.c, E.cpp, E.java, ou E.py}

{\large Timelimit: $1$}

{\tiny Por Márcio Belo}

\end{center}			

O professor Márcio Belo sempre gostou de charadas matemáticas. Um dia ele se deparou com uma estória, pra lá da Estônia. Dado um número primo ímpar $p$, sempre é possível um triângulo de lados inteiros $x$, $y$ e $z$, em ordem crescente de tamanho com as seguintes propriedades: 
\begin{itemize}
    \item A diferença da soma dos lados menores do triângulo pelo lado maior é igual ao primo $p$; 
    \item A diferença da soma das áreas dos quadrados de lados $x$ e $y$ e a área do quadrado de lado $z$ é igual ao primo $p$. 
\end{itemize}
Em matematiquês:

\vspace{3mm}
Dado um primo $p \geq 3$, sempre existem triplas $(x,y,z)$, com $x \leq y \leq z$ que satisfazem: 
\begin{align*}
    x + y - z & = p \\
    x^{2} + y^{2} - z^{2} & = p 
\end{align*}
Quem são essas triplas $(x,y,z)$? Você pode ajudar o professor?

\vspace{0.3cm}
\textbf{Entrada}
Cada entrada contém uma única linha com o número primo $p$, com $(3 \leq p < 2^{16})$. 

\vspace{0.3cm}
\textbf{Saída}
O sistema deve retornar todas as triplas $(x,y,z)$ em ordem crescente de $x$. Uma tripla por linha,  

\begin{table}[h]
    \centering
    \begin{tabular}{|p{7cm}|p{7cm}|}
        \hline
        \begin{center}
            \vspace{-0.3cm}
            \textbf{Exemplos de Entrada}
            \vspace{-0.3cm}
        \end{center} 
         &  
        \begin{center}
            \vspace{-0.3cm}
            \textbf{Exemplos de Saída}
            \vspace{-0.3cm}
        \end{center} \\ \hline
         \texttt{3} & \texttt{(4,6,7)} \\  
         \hline
         \texttt{5} & \texttt{(6, 15, 16)} \\ 
          & \texttt{(7, 10, 12)} \\ 
         \hline
         \texttt{13} & \texttt{(14, 91, 92)} \\
         & \texttt{(15, 52, 54)} \\
         & \texttt{(16, 39, 42)} \\
         & \texttt{(19, 26, 32)} \\ 
         \hline
         \texttt{65267} & \texttt{(65268,2129923278,2129923279)} \\ 
          & \texttt{(97900,130534,163167)} \\ 
         \hline
    \end{tabular}
    \caption{Exemplos de entradas e saídas}

\end{table}

\newpage
\begin{center}
\noindent
\lstinputlisting{abobora_25/E/E5.cpp}
\end{center}